\scanpage{156}\mspage{41}
$\unsure{\mathfrak{M}|\mathscr{M}}(V)\cong\mathfrak{M}(V^{*})$, we infer from the proposition i) that $\widetilde{\Phi}$ induces an isomorphism $\mathfrak{M}(\widetilde{V})\cong\mathfrak{M}(V^{*})$.

The inverse mapping $\widetilde{\Phi}^{-1}$ has possibly a finite number of fundamental points and branch curves and, except for these \uline{singular points}, $\widetilde{\Phi}^{-1}$ is locally bi-regular. Hence, for any general point $\mathfrak{p}^{*}$ on $V^{*}$, the inverse image $\widetilde{\Phi}^{-1}(\mathfrak{p}^{*})$ consists of a finite number of points and the number $m$ of these points is independent of $\mathfrak{p}^{*}$. Now we show that $m=1$. Let $D_h^{*}$ be the section of $V^{*}$ cut out by a general hypersurface of order $h$ containing no fundamental point of $\widetilde{\Phi}^{-1}$ and set $D_h=\widetilde{\Phi}^{-1}(D_h^{*})$. Then, since $\widetilde{\Phi}$ induces the isomorphism $\mathfrak{M}(\widetilde{V})\cong\mathfrak{M}(V^{*})$, \uline{the complete linear system} $\lvert D_h\rvert$ \uline{on} $\widetilde{V}$ \uline{and the complete linear system} $\lvert D_h^{*}\rvert$ \uline{on} $V^{*}$ \uline{have the same diemsion}. For sufficiently large $h$, we infer from the theorem of Riemann-Roch (2.10) and Theorem 2.4 that
\[
\dim\lvert D_h^{*}\rvert=\frac{1}{2}(D_h^{*2})-\frac{1}{2}(K^{*}D_h^{*})+\chi(V^{*})-1,
\]
where $K^{*}$ is the canonical bundle on $V^{*}$. Setting $2a=(D_1^{*2})$, $2b^{*}=(K^{*}D_1^{*})$, we get therefore
\[
\dim\lvert D_h^{*}\rvert=ah^2-b^{*}h+\chi(V^{*})-1.
\]
On the other hand, we get, using (2.10),
\[
\dim\lvert D_h\rvert\geqq\frac{1}{2}(D_h^{2})-\frac{1}{2}(KD_h)+\chi(\widetilde{V})-1-i(D_h).
\]
The general member $D_h^{*}$ of $\lvert D_h^{*}\rvert$ is an irreducible non-singular curve on $V^{*}$. Let $D_h^{*\prime}$ be another general member of $\lvert D_h^{*}\rvert$ and
% archive sheet number: 155
